\documentclass[a4paper,12pt]{article}
\usepackage[T2A]{fontenc}
\usepackage[utf8]{inputenc}
\usepackage[russian]{babel}
\usepackage{amsmath,amssymb,mathtools}
\usepackage{helvet}
\renewcommand{\familydefault}{\sfdefault}
\usepackage{icomma}
\usepackage{geometry}
\geometry{left=23mm,right=23mm,top=22mm,bottom=23mm,headheight=24pt,headsep=7mm}
\usepackage{microtype}
\usepackage{tikz}
\usetikzlibrary{arrows.meta,positioning,shapes.geometric,calc}
\usepackage{tabularx,booktabs}
\usepackage{enumitem}
\usepackage{fancyhdr}
\usepackage{xcolor}
\usepackage{setspace}
\usepackage{parskip}
\usepackage{needspace}
\definecolor{accent}{HTML}{008080}
\definecolor{darkaccent}{HTML}{005858}
\definecolor{textgray}{HTML}{313B3B}
\color{textgray}
\setstretch{1.15}
\setlength{\parindent}{0pt}
\setlength{\parskip}{0.48em}
\setlist[itemize]{leftmargin=1.55em,itemsep=0.15em,topsep=0.3em}
\setlist[enumerate]{leftmargin=1.8em,itemsep=0.2em,topsep=0.3em}
\pagestyle{fancy}
\fancyhf{}
\fancyhead[L]{\footnotesize Метод интервалов и рациональные неравенства}
\fancyhead[R]{\footnotesize Ярослав Верещагин}
\fancyfoot[C]{\footnotesize\thepage}
\renewcommand{\headrulewidth}{0.45pt}
\renewcommand{\headrule}{\hbox to\headwidth{\color{accent}\leaders\hrule height \headrulewidth\hfill}}
\newcommand{\sectiontitle}[1]{\needspace{4\baselineskip}\vspace{0.6em}{\large\color{darkaccent}\bfseries #1}\par\vspace{0.15em}}
\newcommand{\minititle}[1]{\needspace{3\baselineskip}{\color{darkaccent}\bfseries #1}\par}
\newcommand{\result}[1]{\par{\color{darkaccent}\bfseries Ответ:} $#1$.}
\begin{document}
\begin{titlepage}
  \thispagestyle{empty}
  \centering
  \vspace*{26mm}
  {\large\color{darkaccent} МАТЕМАТИКА\par}
  \vspace{25mm}
  {\LARGE Чек-лист\par}
  \vspace{12mm}
  {\LARGE\color{accent} Метод интервалов\par}
  \vspace{2mm}
  {\Large\color{accent} и рациональные неравенства\par}
  \vspace{22mm}
  {\normalsize Подготовка к ОГЭ по математике\par}
  \vfill
  {\large Ярослав Верещагин\par}
  \vspace{12mm}
  {\normalsize 26 октября 2026 г.\par}
  \vspace{19mm}
  {\small Лёвин Артём Александрович\par}
  {\small эксклюзивно для Ярослава Верещагина\par}
  \vspace{20mm}
  \begin{tikzpicture}
  \draw[accent,line width=0.9pt] (-5,0) .. controls (-3.7,0.65) and (-1.8,-0.55) .. (0,0)
                              .. controls (1.65,0.5) and (3.45,-0.52) .. (5,0);
  \end{tikzpicture}
  \vspace*{7mm}
\end{titlepage}
\setcounter{page}{1}

\sectiontitle{1. Что показывает метод интервалов}
Решая рациональное неравенство, мы определяем, \emph{где выражение положительно, отрицательно или равно нулю}. Для дроби
\[
 F(x)=\frac{P(x)}{Q(x)},\qquad Q(x)\ne0,
\]
важны значения $x$, при которых числитель $P(x)$ или знаменатель $Q(x)$ обращается в нуль. Эти значения разделяют числовую ось на промежутки постоянного знака.

\minititle{Подготовка выражения}
Приведите неравенство к виду $F(x)>0$, $F(x)<0$, $F(x)\ge0$ или $F(x)\le0$. При необходимости раскройте скобки, вынесите множители и разложите числитель и знаменатель на множители. \textbf{Уравнения $P(x)=0$ и $Q(x)=0$ решайте со знаком равенства.}

\sectiontitle{2. Критические точки и скобки в ответе}
\begin{center}
\renewcommand{\arraystretch}{1.27}
\begin{tabularx}{\linewidth}{@{}p{0.29\linewidth}X X@{}}
\toprule
\textbf{Точка} & \textbf{Строгое $>,<$} & \textbf{Нестрогое $\ge,\le$} \\
\midrule
Нуль числителя & Исключается & Включается, если знаменатель не равен нулю \\
Нуль знаменателя & Исключается & Всегда исключается \\
\bottomrule
\end{tabularx}
\end{center}
Выколотая точка означает круглую скобку; включённая точка --- квадратную. У $-\infty$ и $+\infty$ скобки всегда круглые.

\textbf{Важно:} если одно число обнуляет и числитель, и знаменатель, \emph{его обязательно исключают}. Сокращение общей скобки не возвращает точку в область допустимых значений.

\sectiontitle{3. Как расставлять знаки}
\begin{enumerate}
\item Найдите все действительные нули числителя и знаменателя; расположите их по возрастанию.
\item Возьмите пробную точку \emph{внутри} одного промежутка, где выражение определено, и установите знак.
\item Переходите к соседним промежуткам, учитывая кратность каждого корня.
\item Оставьте только промежутки с нужным знаком и правильно оформите границы.
\end{enumerate}
Число для пробной подстановки выбирайте удобное; $x=0$ подходит только тогда, когда оно лежит внутри проверяемого промежутка и выражение в нём определено.

\sectiontitle{4. Чётная и нечётная кратность}
\begin{itemize}
\item Через корень \textbf{нечётной} кратности знак произведения или дроби меняется.
\item Через корень \textbf{чётной} кратности знак сохраняется.
\end{itemize}
Например, $(x-1)^2=(x-1)(x-1)$: при переходе через $1$ знак такого множителя остаётся положительным с обеих сторон, а в самой точке равен нулю.

\minititle{Разбор примера}
Решим
\[
 \frac{(x-1)^2(x+2)}{x-3}<0.
\]
Критические точки: $-2$ --- нуль числителя; $1$ --- нуль числителя кратности $2$; $3$ --- нуль знаменателя. На промежутке $x>3$ выражение положительно. Переходя влево через $3$, получаем минус; через $1$ знак \emph{не меняется}; через $-2$ снова получаем плюс.
\result{(-2;1)\cup(1;3)}

\sectiontitle{5. Два разных корня и один повторяющийся корень}
\[
 x^2-9=(x-3)(x+3).
\]
Точки $-3$ и $3$ --- \textbf{разные простые корни}. При переходе через каждую из них знак меняется. В выражении $(x-3)^2$ корень $3$ встречается дважды, поэтому знак не меняется. Если $x^2+4=0$, действительных корней нет: квадрат действительного числа не бывает отрицательным.

\minititle{Не путайте коэффициент с корнем}
В выражении $-9(x+3)(x+1)$ число $-9$ --- постоянный множитель. Оно меняет общий знак произведения, но \textbf{не создаёт точку $-9$} на числовой оси.

\sectiontitle{6. Квадрат суммы и разности}
\[
 (x-1)^2=(x-1)(x-1)=x^2-2x+1.
\]
Нельзя возводить в квадрат только отдельные члены разности: $(x-1)^2\ne x^2-1$. Для метода интервалов часто полезно \emph{сохранить} квадрат скобки, чтобы сразу видеть двойной корень.

\sectiontitle{7. Один корень в числителе и знаменателе}
Рассмотрим выражение
\[
 \frac{(x-1)^2}{(x-1)(x-4)},\qquad x\ne1,\quad x\ne4.
\]
При допустимых $x$ оно равно $\dfrac{x-1}{x-4}$, но точка $x=1$ \textbf{остаётся выколотой}. Сначала запишите ограничения исходной дроби, затем сокращайте.

\sectiontitle{8. Замена $t=x^2$}
Если встречаются степени $x^4$ и $x^2$ либо повторяются выражения с $x^2$, может помочь замена
\[
 t=x^2,\qquad t\ge0.
\]
Теперь изучайте знаки по \textbf{оси $t$}. Отрицательные значения $t$ не рассматривайте: они невозможны при действительном $x$.

\minititle{Разбор примера}
Решим
\[
 \frac{(x^2+2)(x^2-4)}{x^2+1}>0.
\]
После замены имеем
\[
 \frac{(t+2)(t-4)}{t+1}>0,\qquad t\ge0.
\]
При $t\ge0$ множители $t+2$ и $t+1$ положительны. Поэтому знак дроби определяется $t-4$, откуда $t>4$. Возвращаемся к $x$:
\[
 x^2>4\quad\Longleftrightarrow\quad x<-2\ \text{или}\ x>2.
\]
\result{(-\infty;-2)\cup(2;+\infty)}

\sectiontitle{9. Самопроверка перед записью ответа}
\begin{itemize}
\item Указаны только нули \emph{числителя и знаменателя}, а не любые числа из выражения.
\item У двойных корней знак сохраняется; у простых --- меняется.
\item При строгом знаке все нули исключены; нули знаменателя исключены всегда.
\item После сокращения сохранены ограничения исходного выражения.
\item После замены $t=x^2$ учтено $t\ge0$ и выполнен возврат к $x$.
\end{itemize}

\clearpage
\thispagestyle{fancy}
\begin{center}
{\Large\color{darkaccent}\bfseries Алгоритм решения рационального неравенства}
\end{center}
\vspace{4mm}
\begin{center}
\begin{tikzpicture}[>=Stealth,
  node distance=8mm and 12mm,
  start/.style={ellipse,draw=darkaccent,thick,minimum width=32mm,minimum height=9mm,align=center,font=\small},
  act/.style={rounded corners=2pt,rectangle,draw=accent,thick,text width=82mm,minimum height=11mm,align=center,inner sep=5pt,font=\small},
  branch/.style={rounded corners=2pt,rectangle,draw=accent,thick,text width=42mm,minimum height=16mm,align=center,inner sep=5pt,font=\small},
  check/.style={diamond,draw=darkaccent,thick,aspect=2.5,text width=40mm,inner sep=1pt,align=center,font=\small},
  arrow/.style={->,thick,draw=textgray},
  every node/.style={text=textgray}
]
\node[start] (a) {Начало};
\node[check,below=7mm of a] (b) {Нужна замена $t=x^2$?};
\node[branch,below left=8mm and 10mm of b] (c) {$t=x^2$, $t\ge0$\\ось переменной $t$};
\node[branch,below right=8mm and 10mm of b] (d) {Оставить $x$\\ось переменной $x$};
\coordinate (join) at ($(c.south)!0.5!(d.south)+(0,-5mm)$);
\node[act,below=4mm of join] (e) {Привести к виду $F>0$, $F<0$, $F\ge0$ или $F\le0$; разложить на множители};
\node[act,below=6mm of e] (f) {Найти нули числителя и знаменателя; отметить запрещённые точки};
\node[act,below=6mm of f] (g) {Расположить точки на оси; проверить знаки и кратности корней};
\node[act,below=6mm of g] (h) {Выбрать промежутки нужного знака; определить круглые и квадратные скобки};
\node[check,below=7mm of h] (i) {Была замена $t=x^2$?};
\node[branch,below left=8mm and 10mm of i] (j) {Вернуться к $x$: решить неравенство для $x^2$};
\node[branch,below right=8mm and 10mm of i] (k) {Записать итоговые промежутки для $x$};
\node[start,below=8mm of $(j.south)!0.5!(k.south)$] (l) {Ответ};
\draw[arrow] (a)--(b);
\draw[arrow] (b.south west) -| node[pos=0.18,above left,font=\footnotesize]{да} (c.north);
\draw[arrow] (b.south east) -| node[pos=0.18,above right,font=\footnotesize]{нет} (d.north);
\draw[arrow] (c.south) |- (join);
\draw[arrow] (d.south) |- (join);
\draw[arrow] (join)--(e);
\draw[arrow] (e)--(f);
\draw[arrow] (f)--(g);
\draw[arrow] (g)--(h);
\draw[arrow] (h)--(i);
\draw[arrow] (i.south west) -| node[pos=0.18,above left,font=\footnotesize]{да} (j.north);
\draw[arrow] (i.south east) -| node[pos=0.18,above right,font=\footnotesize]{нет} (k.north);
\draw[arrow] (j.south) |- (l.west);
\draw[arrow] (k.south) |- (l.east);
\end{tikzpicture}
\end{center}

\clearpage
\sectiontitle{Тренировочный блок: Путь А}
Решите десять неравенств. На числовой оси отметьте нули числителя и знаменателя; особое внимание уделите кратным корням и недопустимым значениям.

\minititle{Базовый уровень}
\begin{enumerate}[label=\arabic*.,itemsep=0.62em]
\item $\displaystyle\frac{x+2}{x-4}\ge0.$
\item $\displaystyle\frac{(x-1)(x-4)}{x+2}>0.$
\item $\displaystyle\frac{-9(x+3)(x+1)}{x-2}>0.$
\end{enumerate}
\minititle{Уровень выше среднего}
\begin{enumerate}[label=\arabic*.,start=4,itemsep=0.62em]
\item $\displaystyle\frac{(x-1)^2}{x+2}\ge0.$
\item $\displaystyle\frac{(x-4)^2(x+1)}{x-2}>0.$
\item $\displaystyle\frac{(x-1)^2(x+2)}{x-3}<0.$
\item $\displaystyle\frac{x^2-9}{x+1}\le0.$
\end{enumerate}
\minititle{Сложный уровень}
\begin{enumerate}[label=\arabic*.,start=8,itemsep=0.62em]
\item $\displaystyle\frac{(x-1)^2(x+2)}{(x-1)(x-4)}\ge0.$
\item $\displaystyle\frac{x^2-4}{x^2-9}>0.$
\item $\displaystyle\frac{(x^2+2)(x^2-4)}{x^2+1}>0.$
\end{enumerate}

\clearpage
\sectiontitle{Ответы к тренировочному блоку}
\vspace{0.35em}
\renewcommand{\arraystretch}{1.55}
\begin{tabularx}{\linewidth}{@{}p{13mm}X@{}}
\toprule
\textbf{№} & \textbf{Ответ}\\
\midrule
1 & $(-\infty;-2]\cup(4;+\infty)$ \\
2 & $(-2;1)\cup(4;+\infty)$ \\
3 & $(-\infty;-3)\cup(-1;2)$ \\
4 & $(-2;+\infty)$ \\
5 & $(-\infty;-1)\cup(2;4)\cup(4;+\infty)$ \\
6 & $(-2;1)\cup(1;3)$ \\
7 & $(-\infty;-3]\cup(-1;3]$ \\
8 & $[-2;1)\cup(4;+\infty)$ \\
9 & $(-\infty;-3)\cup(-2;2)\cup(3;+\infty)$ \\
10 & $(-\infty;-2)\cup(2;+\infty)$ \\
\bottomrule
\end{tabularx}

\vspace{10mm}
\sectiontitle{Итоговая памятка}
\begin{tabularx}{\linewidth}{@{}p{0.38\linewidth}X@{}}
\toprule
\textbf{Признак} & \textbf{Что делать}\\
\midrule
Нуль числителя & Проверить строгость знака неравенства.\\
Нуль знаменателя & Всегда исключить точку.\\
Корень чётной кратности & Сохранить знак при переходе через него.\\
Разные корни $-a$ и $a$ & Отметить оба корня; проверить их кратности.\\
Квадрат скобки $(x-a)^2$ & Рассматривать как два одинаковых множителя.\\
Замена $t=x^2$ & Учесть $t\ge0$ и выполнить обратную замену.\\
\bottomrule
\end{tabularx}
\vfill
\begin{center}
{\footnotesize\color{darkaccent} Лёвин Артём Александрович --- эксклюзивно для Ярослава Верещагина}
\end{center}
\end{document}
